% !TeX root = ../xdupgthesis_template_lc.tex

% 第一章 绪论（7页控制版）
% 说明：本版在上一稿基础上进行压缩，重点收缩了1.1背景铺陈、1.2.1技术对比和各小节重复总结段。
% 参考文献暂未插入，后续统一补充。

\section{研究背景及研究意义}
\label{sec:intro_background}

随着城市化进程的持续推进和机动车保有量的不断增长，道路交通运行压力日益加大，交通拥堵、停车资源紧张以及道路异常事件处置不及时等问题逐渐成为制约城市交通高效运行的重要因素。车流规模、车型结构、停车行为以及道路占用状态等信息，是交通运行评估、信号控制优化、道路资源配置和异常事件监测的重要依据。因此，如何以较低成本、较高可靠性和较强环境适应性获取道路侧车辆信息，已成为智慧交通建设中的基础性问题。

智慧交通系统通过融合传感器网络、无线通信、边缘计算和智能分析等技术，实现对道路、车辆及交通设施的动态感知与协同管理，是提升道路运行效率和交通治理能力的重要技术途径。对于道路侧长期在线部署的感知系统而言，除检测结果的准确性和稳定性外，还必须满足功耗低、成本可控、维护方便以及适合规模化铺设等工程要求。因此，面向真实道路环境的感知方法不仅要关注识别性能，还要兼顾长期运行能力、复杂工况适应性和系统可部署性。

目前，道路车辆感知手段主要包括视频、毫米波雷达、感应线圈以及地磁传感器等。视频和雷达能够提供较丰富的车辆外观、位置与轨迹信息，但对光照、天气和安装条件较为敏感，且存在成本较高、带宽需求较大和维护复杂等问题；感应线圈技术成熟，但施工和后期维护代价较高。相比之下，地磁传感器通过感知车辆铁磁性结构对局部磁场的扰动实现检测，具有体积小、功耗低、成本较低和便于地埋安装等特点，更适合道路侧低成本、长期在线和分布式部署场景。

对于单地磁节点道路车辆感知系统而言，传感器的观测对象通常是连续采集的三轴地磁序列。车辆通过传感器邻域时，会使局部磁场发生变化，从而在时间轴上形成具有一定幅值、持续时间和形态特征的扰动波形。通过对这些波形进行检测、分割与分析，可以进一步实现车辆通过检测、车型识别以及停车与占用状态判定等任务。与多传感器融合方案相比，单地磁方案在硬件构成和系统部署上更加简洁，特别适用于道路侧节点数量较多、供电条件受限且通信资源紧张的应用场景。然而，单地磁节点可获取的信息维度有限，无法直接获得车辆外观、尺寸和轨迹等显性特征，其有效信息主要隐含在扰动波形的时间结构和幅值变化之中。因此，如何在单传感器条件下充分挖掘磁场信号中的有效信息，是该类系统研究中的核心问题。

从实际应用条件来看，单地磁车辆感知仍面临多方面挑战。首先，环境磁场会随时间发生缓慢变化，传感器自身也可能存在温漂和长期漂移，道路周边偶发性电磁干扰还会引入异常波动，这些因素都会影响车辆事件边界的准确定位以及后续判定结果的稳定性。其次，同一车型在不同速度下通过传感器时，其扰动波形往往会在时间轴上表现出明显的压缩与拉伸，导致样本之间持续时间不一致、局部峰谷位置错位，从而增加车型分类难度。再次，道路侧节点通常采用低功耗硬件平台和无线通信方式运行，端侧算法不仅要具备较好的识别性能，还必须兼顾运算复杂度、存储资源占用和通信开销。最后，在开放道路环境中开展停车与占用检测时，连续车流、邻道影响、慢漂移和伪稳态现象会使简单阈值或固定参考方法难以保持长期稳定，这对停车确认、占用维持和占用释放等状态边界的判定提出了更高要求。

基于上述背景，围绕单地磁节点开展道路车辆感知研究具有较强的理论意义和工程价值。从理论层面看，该问题涉及弱特征条件下的时间序列建模、复杂环境中的状态判定以及单传感器约束下的信息表达与利用等内容，具有一定研究意义；从工程层面看，若能够在单节点、低功耗和低采样率约束下稳定实现车辆检测、车型分类和停车占用识别，将有助于提升道路侧感知系统的部署效率和应用范围。本文面向道路地埋式单地磁节点场景，以连续三轴地磁观测数据为基础，在车辆事件检测与分割的前提下，重点研究车型三分类方法和停车与占用检测方法，力图在不增加传感器数量的条件下进一步挖掘单地磁节点的应用潜力。


\section{国内外研究现状}
\label{sec:related_work}

围绕道路车辆感知问题，国内外已经开展了大量研究工作。总体来看，道路侧车辆信息获取通常遵循"连续观测、事件检测、任务识别"的处理链路，即首先从连续传感序列中完成车辆事件检测与分割，再在事件级数据基础上开展车型识别、停车检测或占用状态判定等任务。结合本文研究内容，下面主要从道路车辆感知技术及地磁车辆检测方法、基于单地磁传感器的车型分类方法以及基于单地磁传感器的停车与占用检测方法三个方面，对相关研究现状进行综述。

\subsection{道路车辆感知技术及地磁车辆检测方法研究现状}
\label{sec:related_sensor_detection}

道路车辆检测是交通信息采集体系中的基础环节，其检测精度和运行稳定性直接影响后续交通流统计、车型识别、停车管理以及异常事件监测等应用效果。现有道路车辆感知技术主要包括视频、毫米波雷达、感应线圈以及地磁传感器等。视频检测信息丰富，但对光照、天气和遮挡较为敏感；毫米波雷达具有较好的环境适应性，但成本和部署复杂度相对较高；感应线圈检测效果稳定，但施工和维护代价较大。相比之下，地磁传感器具有体积小、功耗低、成本较低和便于隐蔽安装等特点，因此在道路侧长期在线部署场景中具有较好的应用前景。

围绕地磁车辆检测，现有研究主要从磁场偏移、背景基线更新、差分特征提取和状态机约束等方面展开。较早的方法多采用阈值和基线更新实现车辆到达与离开判断，后续研究则进一步引入一阶差分、概率判决和时序状态约束，以提高复杂环境下车辆事件检测的稳定性。在低采样率和低功耗条件下，如何兼顾检测精度与算法复杂度，也是相关研究关注的重要问题。

总体来看，基于地磁传感器的车辆检测研究已经形成较为清晰的技术基础。对于本文而言，车辆事件检测与分割是后续车型分类和停车占用判定的前提，因此需要在绪论中进行必要交代，但本文的主要研究重点并不在基础通过检测本身，而在其上层任务的方法设计与性能提升。

\subsection{基于单地磁传感器的车型分类方法研究现状}
\label{sec:related_vehicle_classification}

车型分类是交通流分析中的重要内容。不同类型车辆在道路占有特性、通行能力、能耗水平和安全风险等方面均存在明显差异，因此车型识别结果可为交通规划、收费管理、拥堵治理和道路运行评估提供重要依据。传统车型分类方法主要依赖视频、激光雷达等能够直接获取车辆外观或尺寸信息的检测设备。这类方法识别依据直观，分类边界较为清晰，但其部署维护成本较高，对环境条件的依赖较强，不利于低成本、大范围道路侧部署。

随着磁传感技术的发展，利用单地磁传感器实现车型分类逐渐成为研究热点。由于单地磁节点无法直接获得车辆长度、宽度和外观图像等显性信息，相关研究通常从车辆通过时产生的地磁扰动波形中提取幅值、持续时间、峰谷特征、能量分布、过零点数、统计矩和熵值等人工特征，再结合支持向量机、决策树、随机森林、逻辑回归或 Softmax 分类器完成车型识别。该类方法具有实现简单、可解释性较强、端侧计算开销较低等优点，便于在单节点资源受限场景中落地应用，因此在现有单地磁车型分类研究中占据重要位置。

但是，基于人工特征和轻量分类器的方法也存在较为明显的局限。一方面，不同道路场景、不同安装位置以及不同交通流条件下采集到的波形分布存在差异，手工设计的特征对具体采集条件较为敏感，跨场景泛化能力有限；另一方面，车辆速度变化对单地磁波形影响显著，同一车型在不同速度下经过传感器时，其事件持续时间和局部结构往往存在较大差异，这会直接削弱基于固定时间尺度特征的稳定性。也就是说，即使车辆类别相同，波形在时间轴上的伸缩变化也会使传统特征方法难以始终保持稳定的分类性能。

为了减弱人工特征依赖并提高复杂样本条件下的识别能力，近年来部分研究开始采用一维卷积神经网络、循环神经网络以及其他深度学习模型直接对地磁时序信号进行建模。与传统方法相比，深度模型能够从原始波形中自动提取局部模式和高层特征，在一定程度上提升复杂工况下的分类能力。同时，也有研究尝试通过多传感器融合增强对车辆结构特征的表达，例如利用双磁传感器估计车辆长度或结合其他传感信息提高识别精度。总体来看，深度学习和多传感器路线为车型分类提供了新的思路，但也带来了训练样本需求增加、模型部署复杂度上升以及系统成本提高等问题，这与单地磁节点低成本、易部署的工程目标并不完全一致。

车速变化所引起的波形时间伸缩问题，是影响单地磁车型分类性能的重要因素之一。相同车型在不同速度条件下通过传感器时，磁扰动波形虽然在整体形态上可能保持一定相似性，但其在时间轴上的压缩与拉伸会导致样本之间难以直接对齐，从而影响分类模型对时序结构的利用效率。已有研究通常采用归一化、重采样等方式缓解长度差异，但对于如何在保持工程可实现性的同时更充分地利用波形时序信息，仍然缺乏统一而有效的方法。总体来看，现有单地磁车型分类方法在端侧可部署性与复杂工况鲁棒性之间仍存在权衡，围绕车速变化条件下的时间伸缩问题开展进一步研究具有现实意义。

\subsection{基于单地磁传感器的停车与占用检测方法研究现状}
\label{sec:related_parking_detection}

停车与占用检测是智慧停车管理和道路异常事件监测中的重要研究内容。与车辆通过检测不同，停车检测关注的是车辆由动态经过转变为停留状态的过程，而占用检测则更强调某一位置或区域在一段时间内是否被车辆持续占据。从任务属性上看，停车更接近事件层面的发生与结束，占用更接近状态层面的建立、维持与释放。因此，该类任务不仅需要识别车辆是否停靠，还需要进一步描述占用状态何时形成、持续多久以及何时解除，这对信号边界定位和状态维护能力提出了更高要求。

现有基于磁传感器的停车检测方法中，较为常见的是阈值判决与背景基准更新相结合的技术路线。该类方法通常通过监测磁场偏移量、持续时间以及滑动窗口统计量来判断某一位置是否进入占用状态，并通过基准更新适应环境磁场的缓慢变化。这类方法实现简单、计算量小、易于在资源受限节点上实现，因此在停车位占用监测等相对简单的场景中具有一定实用价值。然而，当环境变化较大或车辆连续经过时，基准更新过程容易受到干扰，导致参考量不再可靠，从而引发误判和漏判。

为了增强停车与占用检测的鲁棒性，研究者进一步引入相关性、方差、极差、均值差等稳定性判据，对停车和占用状态进行辅助识别。通过在一定时间窗口内刻画信号波动程度，可以在一定程度上区分短时扰动与持续停留状态，提高占用检测的可靠性。与此同时，也有研究采用多传感器协同方式，将磁传感器与无线信号、压力传感器、红外传感器等信息进行融合，从而提升系统对复杂工况的适应能力。这类方法在封闭式停车位场景中取得了较好效果，但也不可避免地增加了系统成本和部署复杂度，不完全符合道路侧单地磁节点低成本、易部署的应用目标。

需要指出的是，已有停车与占用检测研究中，相当一部分工作主要面向停车位管理或相对静态的局部场景，而开放道路环境中的停车与占用检测问题更加复杂。首先，道路侧环境磁场可能随时间出现慢变漂移，使得固定基准难以长期保持有效；其次，连续车流会使事件结束后的信号无法迅速回归到稳定无车状态，影响停车确认和占用释放判断；再次，邻道车辆和周边干扰可能形成看似平缓但并不真正稳定的伪稳态信号，这会导致基于简单阈值或单一稳定判据的方法出现较大误差。因此，在道路侧单地磁节点条件下，如何构造更适合复杂工况的动态参考量，并在停车事件与占用状态之间建立更稳定的判定关系，是当前研究中仍需进一步解决的重要问题。

综上所述，现有停车与占用检测方法在轻量实现和复杂工况鲁棒性之间仍然存在矛盾，如何面向慢漂移、连续车流和伪稳态场景构造稳定判定机制，仍是有待解决的问题。


\section{研究内容及章节安排}
\label{sec:intro_content}

\subsection{研究内容}
\label{sec:research_content}

本文面向道路地埋式单地磁节点应用场景，围绕单地磁条件下的道路车辆感知问题展开研究。论文以三轴地磁连续观测数据为基础，在车辆事件检测与分割的前提下，分别针对车型三分类和停车与占用检测两类任务开展研究。

首先，针对连续观测序列中的车辆事件提取问题，本文研究地磁车辆检测的基本原理、事件检测方法以及相应的系统实现方式。该部分工作主要围绕车辆经过传感器时所产生的磁扰动机理展开，结合路侧检测节点、通信链路和平台系统，明确车辆事件检测与数据获取的基础过程。该部分内容虽然不是全文最终希望解决的高层任务，但它决定了后续车型分类和停车占用判定所依赖的输入质量与数据组织方式，因此构成全文研究展开的前提。

其次，针对道路侧单地磁条件下的车型三分类问题，本文以小型车、中型车和大型车为研究对象，构建基于车辆事件片段的分类方法。考虑到道路侧节点在算力、功耗和通信资源方面存在约束，论文首先研究基于统计特征与 Softmax 分类器的轻量化分类方法，以满足端侧部署需求；与此同时，针对车辆速度变化引起的波形时间伸缩现象，进一步研究基于动态时间规整和一维卷积神经网络的离线分类方法，用以分析在时间对齐条件下单地磁波形的可分性及其性能提升空间。通过这两条研究路径，本文试图在工程可实现性与分类性能之间建立较为合理的平衡。

再次，针对道路场景中的停车与占用检测问题，本文研究基于信号稳定点的状态判定方法。考虑到固定参考或简单阈值方法在慢漂移、连续车流和伪稳态条件下容易出现误判，论文从车辆事件结束前后信号由动态扰动向相对稳定状态过渡的过程入手，构造停车前后动态参考量，并结合时序约束与状态维护机制，实现停车确认、占用维持和占用释放的统一描述。该部分工作的重点不在于给出一个仅适用于理想场景的静态规则，而是在单地磁、开放道路和复杂干扰背景下提高状态判定的稳定性与鲁棒性。

总体而言，本文并不是将车辆检测、车型分类和停车占用检测三个问题简单并列，而是在单地磁连续观测基础上形成一条相互衔接的研究链路。与已有仅关注单一任务的研究相比，本文更加重视在单节点、50 Hz 采样和道路侧长期运行约束下对同一类观测数据的多层次利用：在车型分类方面，重点解决车速变化带来的时间伸缩问题；在停车与占用检测方面，重点解决复杂工况下参考量难以稳定维护的问题。本文希望在不增加传感器数量的前提下，进一步挖掘单地磁节点的应用潜力，并为后续工程部署与系统扩展提供方法基础。

\subsection{章节安排}
\label{sec:thesis_organization}

本文共分为五章，各章节内容安排如下。

第一章为绪论。本章主要介绍课题的研究背景及研究意义，分析单地磁道路车辆感知的应用需求，综述相关领域的国内外研究现状，并在此基础上给出本文的研究内容和章节安排。

第二章介绍地磁车辆事件检测基本原理与系统设计。本章从车辆磁扰动形成机理出发，研究连续三轴地磁序列中的车辆事件检测与分割基础，并结合路侧检测节点、通信方式和平台系统，对数据采集与处理流程进行说明，为后续章节研究提供基础支撑。

第三章研究基于单地磁传感器的车型分类方法。本章围绕小型车、中型车和大型车三分类任务，首先构建事件级样本并研究基于统计特征的轻量分类方法，随后进一步分析车速变化对波形时域结构的影响，研究基于动态时间规整和一维卷积神经网络的离线分类方法，并通过实验对不同方法的性能进行比较与分析。

第四章研究基于信号稳定点的道路停车与占用检测方法。本章结合道路场景下停车与占用状态的特点，分析相关信号变化规律，研究基于动态参考的停车与占用判定方法，并在不同复杂工况下对算法的有效性和鲁棒性进行验证与分析。

第五章为总结与展望。本章对全文工作进行总结，归纳本文研究的主要结论，并结合当前工作中仍存在的问题，对后续在模型轻量化、复杂工况适应性提升以及工程部署扩展等方向上的研究进行展望。

